% GENERATED, edit the French unit or the translation table
\mxnewpage\section*{EV1-T1 -- Common test}\label{EV1-T1}
\addcontentsline{toc}{section}{EV1-T1 -- Common test}
\textbf{Practice paper}. Chapters 1 to 4.\par Duration: 2 h. No calculator. Marked out of 20. 14 November 2026.
\subsection*{Part 1 -- Fluency (6 points)}
Twelve questions, 0.5 point each. Only one correct answer; no justification required. An incorrect or missing answer scores 0, with no loss of points.
\begin{enumerate}
\item Calculate $3+2\times3$.\par A. $12$\quad B. $9$\quad C. $8$\quad D. $15$
\item Calculate $\frac12+\frac13$.\par A. $\frac{5}{3}$\quad B. $\frac{2}{5}$\quad C. $\frac{5}{6}$\quad D. $\frac{1}{6}$
\item Calculate $\dfrac{2^3}{2^{1}}$.\par A. $4$\quad B. $8$\quad C. $2$\quad D. $1$
\item A value of 50 increases by 20 \%. Final value?\par A. $10$\quad B. $40$\quad C. $60$\quad D. $70$
\item A value of 80 is obtained after a 20 \% decrease. Initial value?\par A. $60$\quad B. $64$\quad C. $96$\quad D. $100$
\item A 25 \% increase then a 20 \% decrease: overall coefficient?\par A. $\frac{4}{5}$\quad B. $\frac{21}{20}$\quad C. $\frac{9}{10}$\quad D. $1$
\item The solution of $2x+3=0$ is\par A. $\frac{3}{2}$\quad B. $\frac{-3}{2}$\quad C. $6$\quad D. $-6$
\item The slope between $(0;3)$ and $(2;9)$ is\par A. $2$\quad B. $-3$\quad C. $6$\quad D. $3$
\item The mean of $3, 5, 7$ is\par A. $6$\quad B. $15$\quad C. $3$\quad D. $5$
\item A sector of $90^\circ$ represents what fraction of the total?\par A. $\frac{1}{2}$\quad B. $\frac{1}{4}$\quad C. $\frac{3}{4}$\quad D. $\frac{1}{3}$
\item A probability is $\dfrac{2}{5}$. Its complement is\par A. $\frac{2}{5}$\quad B. $\frac{3}{5}$\quad C. $\frac{5}{2}$\quad D. $\frac{1}{5}$
\item For $f(x)=x^2+3$, $f(-2)$ is\par A. $7$\quad B. $5$\quad C. $12$\quad D. $-1$
\end{enumerate}
\mxnewpage\subsection*{Part 2 -- Independent exercises (14 points)}
Every step must be justified. Exact results are expected unless stated otherwise.
\mxneed{12}
\subsection*{Exercise 1 -- A book collection (5 points)}
A library has 120 books at the end of the reference month, called month 0. At the end of each following month, it adds 18 books to its stock; no book is removed. Let $u_n$ be the number of books available at the end of month $n$, for every natural number $n$.
\begin{enumerate}
\item Determine the type of $(u_n)$, its common difference and its expression in terms of $n$. \hfill(1 pt)
\item Calculate $u_8$, then $u_0+u_1+\cdots+u_8$. Explain what this sum accumulates and why it does not count books that are all distinct. \hfill(2 pts)
\item Determine the first month in which the stock strictly exceeds 300 books. Check the previous month. \hfill(2 pts)
\end{enumerate}
\mxneed{12}
\subsection*{Exercise 2 -- Area and choice of form (5 points)}
A rectangle has a perimeter of 28 m. If one of its sides measures $x$ metres, the other measures $14-x$ metres. Its area is modelled by the function $A$ defined on $[0;14]$ by $A(x)=x(14-x)$, in square metres. The endpoints represent degenerate rectangles.
\begin{enumerate}
\item Expand, then write A in vertex form. \hfill(1 pt)
\item Determine the maximum area and the corresponding value of x without a derivative. \hfill(2 pts)
\item Solve $A(x)\geq40$ on the domain. \hfill(2 pts)
\end{enumerate}
\mxneed{12}
\subsection*{Exercise 3 -- A right angle (4 points)}
The plane is equipped with an orthonormal frame. Given the points $A(0;1)$, $B(3;2)$ and $C(-1;4)$. Lengths are expressed in the unit of the frame and areas in area units.
\begin{enumerate}
\item Calculate the coordinates of the vectors $\vect{AB}$ and $\vect{AC}$. \hfill(1 pt)
\item Prove that these vectors are orthogonal and state the nature of the triangle. \hfill(1 pt)
\item Calculate the length $BC$ then the area of the triangle $ABC$. \hfill(2 pts)
\end{enumerate}
